% SPDX-License-Identifier: CC-BY-4.0
% Copyright (c) 2026 Martin Schröder <info@swedishembedded.com>

\section{Introduction}\label{sec:intro}

A user who wants a model that reasons and writes ``like'' a historical author
has, today, two cheap options. The first is to tell a general model who to be, in
a system prompt. The second is to fine-tune a model on the author's collected
writings, optionally with synthetic dialogues about them. Whether the second option
buys anything over the first, and how one would know, is the question this paper
studies. We study it inside a concrete system, Splinter, which is given a single
sentence of the form ``Learn to think like \emph{X} based on the materials \emph{X}
has written in directory \emph{D}'' and must, without further instruction, produce a
fine-tuned model together with a measurement of what it learned and a decision on
whether to release it.

We chose two test subjects with large public-domain corpora, Thomas Jefferson (about 1.5
million words of letters and works in the training materials) and Samuel Adams (a smaller
corpus of about one megabyte). The policy is Qwen3-8B~\citep{qwen3} used in non-thinking
mode and adapted with LoRA~\citep{hu2021lora}; the judge is a separate Qwen3-14B. Everything
runs on two Tesla P40 cards with 24~GB each, which constrains the work in ways that shape
the methodology (\cref{sec:residency}).

\paragraph{Scope.}
This paper covers one question: whether fine-tuning a small policy model on one
author's own writings makes it think and write like that author on material it has not
seen, beyond what a persona prompt achieves, and how that is measured. In scope are the
construction of the training data from the author's writings (views, tasks, abstention,
replay), the training mechanics of the adapter, the evaluation design (held-out source
families, exam, judge, statistics, retention), the release gate's checks as measurements,
and the defects in any of these. Out of scope, and reported elsewhere when they have
verified findings of their own, are other uses of the system (timelines and longitudinal
learning, predictive gates, learning from usage policy), retrieval of source passages at
answer time, serving performance, and models other than the Qwen3 pair studied here. A
new finding belongs in this paper if it changes what is claimed about, or how one would
measure, an adapter trained on an author's writings.

\paragraph{The central difficulty is measurement, not training.}
It is easy to produce an adapter that loses a few points of training loss and
sounds a little more archaic. It is hard to show that the adapter has changed what the
model \emph{does} on text it has never seen, beyond what a prompt already achieves, at a
cost to general behaviour that is acceptable. The defects that changed the reported
numbers of the first version of the system were defects of measurement: a split that
leaked, a gate that measured the wrong thing, a grader that failed correct answers, an
exam with too few independent units to detect anything. The training defects were found
by reading the training code, not through the exam numbers.

\paragraph{Contributions.}
\begin{enumerate}
\item A complete, autonomous pipeline from one sentence to a released adapter, with a
  family-level held-out rule that is robust to overlapping editions (\cref{sec:system,sec:split}).
\item An evaluation design for small-corpus persona fine-tuning: an exam reserved
  before any generation, a primary comparison fixed in advance against the
  \emph{persona-prompted} base, per-task pairing with a family-clustered bootstrap, Holm
  correction of secondaries, a power simulation for sizing, and a judge with controls and
  deterministic overrides (\cref{sec:eval}). The design is implemented, its power is
  simulated, four pilot runs on a subset of the exam and full runs of two candidates on all of
  it are reported (\cref{sec:frozen,sec:deployed,sec:full}); the other results of \cref{sec:results} come from an earlier,
  simpler exam (\cref{sec:exam-versions}).
\item A catalogue of defects that were found and fixed, each with the measurement that
  exposed it: leakage across editions, a gate that did not test generalisation, a grader bug,
  the retention suite's inability to detect a two-point drop, and eight training-mechanics
  defects, seven of them corrected (\cref{sec:defects}).
\item Results (\cref{sec:results}). On the earlier exams no candidate is
  distinguishable from the persona-prompted base at $\alpha=0.05$. On the first 25 families
  of the reserved exam, the v7 candidate is not distinguishable from it either (+4.5
  points, family-level $p=0.21$), while the candidate trained with voice data is judged
  right on 11.0 points more tasks ($p=0.004$) only when it is also given the persona
  prompt. Two further candidates, run with only the deployed arm, give one more positive
  result (trained with voice data under the corrected recipe: +8.4 points, family-level
  $p=0.018$) and one that is not distinguishable (+4.8 points, $p=0.30$)
  (\cref{sec:deployed}). On the full frozen exam (242 tasks, 38 families) the two positive
  candidates keep their lead over the prompted base under the persona prompt, at a smaller size:
  +9.1 points (family-level $p=0.0033$; $0.013$ after a factor of four for the candidates
  examined on the reserved exam) and +7.0 points ($p=0.022$; $0.087$). The exam contains the pilot, the
  candidates were chosen after it, the part outside the pilot does not resolve the leads, and
  neither adapter exceeds the base or the prompted base without the persona prompt
  (\cref{sec:full}). We give the measured numbers, the bounds on what the exams could
  have shown, and what has not been established.
\end{enumerate}

\paragraph{What this paper does not claim.}
It does not claim that the adapter ``thinks like'' either author. It does not claim that
a fine-tune beats prompting without the prompt: on the earlier exams the candidates' point
estimates are at or above the prompted base but cannot be separated from chance, and on the
reserved exam no candidate asked without the persona prompt exceeds the prompted base (two
candidates were not run without it, and the two that were run on the whole exam fall below
it); the distinguishable results are for candidates that are also given the prompt, chosen after
earlier exams and after the pilot, on an exam that contains the pilot. Numbers
that were not measured are absent or marked as such;
\cref{sec:repro} lists the commands that reproduce each measured number and
\cref{sec:limits} what remains unestablished.
